\pagestyle{plain}
\chapter*{Resumen}

La gestión eficiente de las redes de transporte público presenta desafíos progresivamente mayores debido a la necesidad de integrar grandes volúmenes de datos tanto espaciales como temporales. Ante este escenario, surge la necesidad de evaluar si las arquitecturas SQL tradicionales siguen siendo óptimas, o por el contrario, los paradigmas orientados a grafos ofrecen ventajas significativas.

En este trabajo se han modelado dichas redes geoespaciales en dos sistemas gestores de bases de datos diferentes, PostgreSQL y Neo4J, partiendo del estándar GTFS. Después, se ha desarrollado un conjunto representativo de consultas que responden a necesidades de información reales, tanto de usuarios de la red como de las empresas que ofrecen dichos servicios.

A continuación, se ha trabajado en la representación gráfica, sobre el globo terráqueo, de los resultados generados, utilizando QGIS y Folium, y, en último lugar, se ha llevado a cabo una comparativa entre las prestaciones de ambos sistemas, centrada no solo en la eficiencia computacional, sino también en la facilidad de modelado y expresión de consultas.

Los resultados obtenidos muestran que PostgreSQL ofrece un rendimiento superior en ingesta de datos y agregación masiva de estos, contando además con excelente soporte para información geoespacial. A su vez, Neo4J destaca en consultas que requieren navegación desde un número reducido de nodos iniciales, y ofrece una sintaxis más legible y fácil de mantener, aunque presenta limitaciones en el soporte de bibliotecas.
